\section{2020 Geometry \& Topology}

\subsection{Exam}

\begin{exercise}
(a) Show that $\mathbb{RP}^{2n}$ cannot be the boundary of a compact manifold.

(b) Show that $\mathbb{RP}^{3}$ is the boundary of some compact manifold.
\end{exercise}

\begin{solution}
\textbf{Pedagogical Motivation:} This problem introduces the concept of \textit{cobordism}. A manifold is a boundary if its characteristic numbers vanish. This is a powerful link between topology and algebraic invariants.

\textbf{Formal Proof:}
(a) Recall the theorem of Thom: A compact smooth manifold $X$ is the boundary of a compact smooth manifold $W$ if and only if all its Stiefel-Whitney numbers are zero. 
For the real projective space $\mathbb{RP}^{2n}$, the total Stiefel-Whitney class is $w(\mathbb{RP}^{2n}) = (1+a)^{2n+1}$, where $a \in H^1(\mathbb{RP}^{2n}; \mathbb{Z}_2)$ is the generator.
The top Stiefel-Whitney class is $w_{2n}(\mathbb{RP}^{2n}) = \binom{2n+1}{2n} a^{2n} = (2n+1) a^{2n} = a^{2n} \neq 0$.
The corresponding Stiefel-Whitney number $\langle w_{2n}, [\mathbb{RP}^{2n}] \rangle = 1 \neq 0$.
Since not all Stiefel-Whitney numbers are zero, $\mathbb{RP}^{2n}$ cannot be a boundary.

(b) For $\mathbb{RP}^3$, $w(\mathbb{RP}^3) = (1+a)^4 = 1 + 4a + 6a^2 + 4a^3 = 1$ in $\mathbb{Z}_2$ coefficients.
All $w_i = 0$ for $i > 0$, so all Stiefel-Whitney numbers are zero.
\textbf{Geometric Intuition:} $\mathbb{RP}^3 \cong SO(3)$. Any Lie group is a boundary. Specifically, $SO(3)$ is the boundary of the unit disk bundle of its own tangent bundle (after some framing adjustments) or more simply, $SO(n)$ is a boundary because it admits a null-cobordism through the group of rotations.
\end{solution}

\begin{exercise}
Suppose $M$ is a noncompact, complete $n$-dimensional manifold, and suppose there is an open subset $U \subset M$ and an open set $V \subset \mathbb{R}^{n}$ such that $M \setminus U$ is isomorphic to $\mathbb{R}^{n} \setminus V$. If $\mathrm{Ric} M \geq 0$, show that $M$ is isometric to $\mathbb{R}^{n}$.
\end{exercise}

\begin{solution}
\textbf{Geometric Intuition:} This is a \textit{rigidity} result. In $\text{Ric} \geq 0$ geometry, the volume of balls grows at most as fast as in $\mathbb{R}^n$. If the "end" of the manifold looks exactly like $\mathbb{R}^n$, then the volume growth at infinity is maximal. This "forces" the interior to also be flat.

\textbf{Formal Proof:}
Suppose $M \setminus K$ is isometric to $\mathbb{R}^n \setminus V$.
Then the volume of large balls $B(p, r)$ in $M$ satisfies:
\[ \lim_{r \to \infty} \frac{\text{Vol}(B(p, r))}{r^n} = \omega_n \]
where $\omega_n$ is the volume of the unit ball in $\mathbb{R}^n$.
By the \textit{Bishop-Gromov Volume Comparison Theorem}, for a manifold with $\text{Ric} \geq 0$, the ratio $\frac{\text{Vol}(B(p, r))}{r^n}$ is a non-increasing function of $r$.
Since the limit as $r \to \infty$ is $\omega_n$, and for small $r$ the limit is also $\omega_n$ (standard for any Riemannian manifold), the function must be constant:
\[ \frac{\text{Vol}(B(p, r))}{r^n} = \omega_n \quad \text{for all } r. \]
The equality case of Bishop-Gromov states that if $\text{Vol}(B(p, r)) = \omega_n r^n$ for some $r$, then the ball is isometric to a Euclidean ball. Thus $M$ is isometric to $\mathbb{R}^n$.
\end{solution}

\begin{exercise}
Compute all the homotopy groups of the $n$-torus $T^{n} = S^{1} \times S^{1} \times \cdots \times S^{1}$, $n \geq 2$.
\end{exercise}

\begin{solution}
The universal cover of $T^n$ is $\mathbb{R}^n$. Since $\mathbb{R}^n$ is contractible, all higher homotopy groups of the torus must vanish.
\[ \pi_k(T^n) \cong \pi_k(\mathbb{R}^n) = 0 \quad \text{for } k \geq 2. \]
The fundamental group is $\pi_1(T^n) \cong \pi_1(S^1)^n \cong \mathbb{Z}^n$.
Thus $T^n$ is an $K(\mathbb{Z}^n, 1)$ space.
\end{solution}

\begin{exercise}
Consider the upper half space $\mathbb{H}^{3}$ with metric $g = \frac{dx^{2} + dy^{2} + dz^{2}}{z^{2}}$. Let $P$ be the surface defined by $\{z = x\tan\alpha, z > 0\}$. Compute the mean curvature of $P$.
\end{exercise}

\begin{solution}
\textbf{Geometric Intuition:} In hyperbolic space $\mathbb{H}^3$, "planes" through the origin in the Euclidean sense are actually \textit{geodesic planes} (hyperrectangles) if they are vertical. However, a tilted plane like $z = x \tan \alpha$ is not totally geodesic unless it's vertical ($\alpha = \pi/2$). The mean curvature measures how much it deviates from being a minimal surface.

\textbf{Formal Proof:}
Let $\bar{g} = \phi^2 g_E$ with $\phi = 1/z$. The mean curvature $\bar{H}$ (sum of principal curvatures) is:
\[ \bar{H} = \frac{1}{\phi} (H_E + (n-1) \frac{\partial \ln \phi}{\partial n_E}) \]
Here $H_E = 0$ for a Euclidean plane. $n=3$.
$\nabla \ln \phi = \nabla \ln (1/z) = (0, 0, -1/z)$.
The unit normal to $z - x \tan \alpha = 0$ is $n_E = (-\sin \alpha, 0, \cos \alpha)$.
Then $\frac{\partial \ln \phi}{\partial n_E} = -\frac{\cos \alpha}{z}$.
\[ \bar{H} = z (0 + 2 (-\frac{\cos \alpha}{z})) = -2 \cos \alpha. \]
If defined as the average, $H = -\cos \alpha$.
\end{solution}

\begin{exercise}
Suppose $M$ is a compact 2-dimensional Riemannian manifold without boundary, with positive sectional curvature. Show that any two compact closed geodesics on $M$ must intersect with each other.
\end{exercise}

\begin{solution}
\textbf{Geometric Intuition:} Positive curvature makes geodesics "bend" toward each other. On a sphere (constant $K>0$), any two great circles intersect at two points. This theorem generalizes this to arbitrary $K>0$.

\textbf{Formal Proof:}
This is \textit{Frankel's Theorem}. Suppose two closed geodesics $\gamma_1, \gamma_2$ are disjoint. Since they are compact, there exists a shortest geodesic $\sigma$ of length $L > 0$ connecting them. $\sigma$ must be perpendicular to both $\gamma_1$ and $\gamma_2$.
The second variation of length for $\sigma$ along a parallel vector field $X$ (tangent to the geodesics) is:
\[ L''(0) = \int_0^L (|\dot{X}|^2 - R(X, \dot{\sigma}, \dot{\sigma}, X)) ds \]
Since $X$ is parallel transported, $\dot{X} = 0$. Since $K > 0$, $R(X, \dot{\sigma}, \dot{\sigma}, X) > 0$.
Thus $L''(0) < 0$. This implies that the distance between $\gamma_1$ and $\gamma_2$ can be decreased, contradicting that $\sigma$ was the shortest path. Thus they must intersect.
\end{solution}

\begin{exercise}
Suppose $\Sigma$ is a smooth compact embedded hypersurface without boundary in $\mathbb{R}^{n}$ for $n \geq 3$. Show that $\Sigma$ is orientable.
\end{exercise}

\begin{solution}
\textbf{Geometric Intuition:} An embedded hypersurface in $\mathbb{R}^n$ always has an "inside" and an "outside". This global separation allows us to pick a consistent normal direction everywhere.

\textbf{Formal Proof:}
By the \textit{Jordan-Brouwer Separation Theorem}, $\mathbb{R}^n \setminus \Sigma$ has exactly two components, one of which (the interior $D$) is bounded.
We can define a global transverse vector field by taking the outward unit normal to $D$. 
A hypersurface is orientable if and only if its normal bundle is trivial. 
Since $\mathbb{R}^n$ is orientable, the tangent bundle $T\mathbb{R}^n$ is trivial. On $\Sigma$, we have the decomposition $T\mathbb{R}^n |_\Sigma = T\Sigma \oplus N\Sigma$.
Since $N\Sigma$ has a non-vanishing section (the outward normal), it is trivial.
Thus $w_1(T\Sigma) = w_1(T\mathbb{R}^n|_\Sigma) - w_1(N\Sigma) = 0 - 0 = 0$.
Hence $\Sigma$ is orientable.
\end{solution}
